% Generated draft: requires mathematical and source review.
\documentclass[11pt]{article}
\usepackage[margin=25mm]{geometry}
\usepackage{fontspec}
\setmainfont{Latin Modern Roman}
\usepackage{amsmath,amssymb,mathtools,unicode-math}
\setmathfont{Latin Modern Math}
\usepackage{enumitem,needspace,xurl,hyperref}
\setlength{\parindent}{0pt}
\setlength{\parskip}{5pt}
\setlength{\emergencystretch}{3em}
\newcommand{\N}{\mathbb{N}}
\newcommand{\Z}{\mathbb{Z}}
\newcommand{\Q}{\mathbb{Q}}
\newcommand{\R}{\mathbb{R}}
\newcommand{\C}{\mathbb{C}}
\newcommand{\F}{\mathbb{F}}
\newcommand{\A}{\mathbb{A}}
\newcommand{\PP}{\mathbb{P}}
\newcommand{\E}{\mathbb{E}}
\newcommand{\eps}{\varepsilon}
\newcommand{\dd}{\,\mathrm{d}}
\begin{document}
\section*{UnsolvedMath ID: 1203}
\textbf{Source problem code:} catalogue entry
\subsection{Definitions and mathematical statement}
For $n\ge1$, let
\[
 S^{n+1}=\{x\in\R^{n+2}:\|x\|=1\}
\]
carry its round metric. A smooth hypersurface immersion $f:M^n\to S^{n+1}$ is minimal if its mean-curvature vector vanishes everywhere, and complete if the induced Riemannian metric on $M$ is complete. The registry asks for a classification of complete minimal hypersurfaces in such spheres for all $n$. It does not specify immersion versus embedding, connectedness, allowance of boundary or singularities, compactness, topology, or the equivalence relation in this classification. Even a restriction to smooth connected boundaryless immersions, with equivalence by reparametrization and ambient isometry, would be a choice requiring confirmation rather than a complete formulation supplied by the record.

The supplied literature triage distinguishes this broad request from Chern's named spherical Bernstein rigidity question. A precise smooth version of the latter, consistent with the supplied embedded and topological scope, asks dimension by dimension: if $M^n$ is a smooth manifold homeomorphic to $S^n$ and $f:M\to S^{n+1}$ is a smooth minimal embedding, must there exist a unit vector $v\in\R^{n+2}$ such that
\[
 f(M)=\{x\in S^{n+1}:\langle x,v\rangle=0\}?
\]
The right-hand side is a totally geodesic equatorial $n$-sphere. Here $M$ is compact and the induced metric is complete; completeness alone does not replace embeddedness and spherical topology. This rigidity question is a different formulation from classifying every complete minimal hypersurface.

The record's literature assessment checked 17 August 2026 reports that Wang, Wang and Zhou constructed a non-equatorial embedded minimal hypersphere in $S^4$ in 2026, giving a negative instance of this rigidity question for $n=3$. It also records known non-equatorial minimal hypersurfaces of other topologies and isoparametric types, and reports no all-dimension classification of the unrestricted class. It labels the situation partial progress, while the registry's main label is open. These are supplied dated assessments, not newly verified claims. The background compares the question with the Euclidean Bernstein problem and mentions dimension-specific progress; its unqualified ``solved by 1968'' gloss does not supply the dimension-dependent Euclidean theorem, which is not asserted here.
\subsection{Short English statement}
The stored question asks broadly to classify complete minimal hypersurfaces in round spheres, without enough scope to specify a classification theorem. Chern's narrower named question concerns embedded minimal hyperspheres; the supplied 2026 assessment reports a non-equatorial example in $S^4$.
\subsection{Sources}
\begin{itemize}
\item Local UnsolvedMath registry, permanent ID $1203$, source code \texttt{catalogue entry}, ``The Spherical Bernstein Problem,'' with its background and literature triage checked 17 August 2026.
\item S.-S. Chern, \emph{Differential geometry: its past and its future}, Proceedings of the International Congress of Mathematicians, Nice 1970, Problem VI, cited in the record for the embedded topological-sphere scope: \url{https://people.math.harvard.edu/~hirolee/pdfs/2014-fall-230a-icm1970-chern-differential-geometry.pdf}.
\item Tongrui Wang, Zhichao Wang and Xin Zhou, \emph{Equivariant min-max theory and the spherical Bernstein problem in $S^4$}, arXiv:2602.03984 (2026), cited for the non-equatorial embedded minimal hypersphere: \url{https://arxiv.org/abs/2602.03984}.
\end{itemize}
\end{document}
